\section{Execution Protocol and Measurement Instrumentation}
\label{sec:execution_protocol}

This section details the physical testbed environment, cluster isolation safeguards, and automated instrumentation protocol used to execute the paired experimental campaigns. We describe how hardware parity is enforced across nodes and how authoritative timestamps and resource vectors are captured to decompose platform actuation delay from observability diagnostic latency.

\subsection{Testbed Isolation and Hardware Parity}
All experiments are executed on an isolated Kubernetes cluster (K3s/K3d) with hard per-node Linux \texttt{cgroup} resource quotas ($2.0\,\text{vCPUs}$ and $4096\,\text{MiB}$ RAM per worker node, $12\,\text{GiB}$ total cluster RAM) running the scoped CNCF OpenTelemetry Astronomy Shop under continuous Locust synthetic user traffic and monitored by the PANOPTES observability stack (OpenTelemetry Collector, Prometheus, Grafana Loki, and Grafana Tempo). To prevent cross-platform interference between the two chaos control planes, Chaos Mesh (deployed in the \texttt{chaos-testing} namespace) and LitmusChaos (deployed in the \texttt{litmus} namespace) are evaluated in dedicated, clean cluster epochs with identical baseline states.

\subsection{Automated Timing and Overhead Instrumentation}
For each trial $i \in \{1, \dots, 30\}$ under each platform treatment, an automated Python orchestration harness records four authoritative UTC timestamps and resource vectors:
\begin{enumerate}
    \item \textbf{CRD Submission Timestamp ($T_{\text{CRD}}$):} Captured immediately when the Kubernetes API Server accepts the Chaos Mesh CRD or LitmusChaos \texttt{ChaosEngine} manifest.
    \item \textbf{Kernel Fault Activation Timestamp ($T_{\text{effect}}$):} Captured by monitoring the exact moment the kernel-level perturbation takes effect on the target pod (i.e., \texttt{tc qdisc} netem rule insertion for network delay, container \texttt{SIGKILL}/deletion event for pod termination, or \texttt{stress-ng} process attachment inside the target container's \texttt{cgroup}). The \textbf{Actuation Latency} is computed as:
    \begin{equation}
        \Delta T_{\text{actuate}} = T_{\text{effect}} - T_{\text{CRD}}
    \end{equation}
    \item \textbf{Diagnostic Convergence Timestamp ($T_{\text{diag}}$):} Captured when the blind PANOPTES multimodal RCA engine (polling Prometheus, Loki, and Tempo HTTP APIs every $4\,\text{s}$) isolates the root-cause microservice. This allows decomposing total diagnostic time ($MTTDiag_{\text{total}} = T_{\text{diag}} - T_{\text{CRD}}$) into platform actuation delay plus net observability diagnostic latency:
    \begin{equation}
        MTTDiag_{\text{total}} = \Delta T_{\text{actuate}} + \underbrace{(T_{\text{diag}} - T_{\text{effect}})}_{MTTDiag_{\text{net}}}
    \end{equation}
    \item \textbf{Teardown Completion Timestamp ($\Delta T_{\text{teardown}}$):} Elapsed wall-clock time between issuing CRD deletion/cancellation and complete restoration of the target container's network/cgroup state (including termination of Litmus \texttt{go-runner}/\texttt{chaos-runner} pods or removal of Chaos Mesh finalizers).
\end{enumerate}

For resource overhead evaluation (RQ02), steady-state operator footprint is sampled across $N=30$ observations at $\Delta t = 16\,\text{s}$ intervals (exceeding the $15\,\text{s}$ Kubelet \texttt{metrics-server} window \cite{hurlbert_1984}), while active-campaign overhead is sampled continuously across the \texttt{chaos-testing} and \texttt{litmus}/\texttt{otel-demo} namespaces during the $N=30$ fault runs.
